\chapter{Automatic speech recognition}

\section{RNN-T}

An RNN-T learns a joint probability of the next token conditioned on the text and acoustic history $P_t(y_{u+1}|\vy_{1:u},\vx_{1:t})$. The probability iterates over the acoustic dimension $t$, and each time the joint emits a token, $u$ increases by 1. The joint is typically implemented by representing the contexts $\vy$ and $\vx$ with two unimodal encoders e.g. a text predictor and an acoustic encoder
\[
P_t(y_{u+1}|\vy_{1:u},\vx_{1:t}) = \vj(\vp(\vy_{1:u}),\ve(\vx_{1:t})).
\]
There are many alignments between $\vy$ and $\vx$, RNN-T solves this by using dynamic programming to optimise over all possible alignments.

\section{Connection to coupled dynamical systems}
Our problem stems from trying to describe the nature of $\vy$ and $\vx$.
\[
y_u=? \qquad x_t=?
\]
First we know that $x_t$ is produced from speech production, and for the purposes of ASR we can leave a mystery. Second, $y$ depends on itself and $x$ through $P_t$, from here we can imagine an equation that could satisfy dependence
\[
y_{u+1}=f(\vy_{1:u}) + g(\vx_{1:t}) + c.
\]
One form of coupling is specifying the rate of change of $u$ w.r.t. $t$
\[
\frac{\dd u}{\dd t}=?
\]
This works nicely as we can specify that a token is emitted every $\Delta u = 1$, and knowing how $u$ depends on $t$ supports this.

For a streaming model, consider a hidden state $\vz_t$ that emits tokens as $p(y_{u+1}|\vz_t)$ when $\Delta u = 1$. The progression of $u$ is slower than $t$, and could be optimised with a neural network. If this is the case, the chain rule becomes useful.
\[
\dd \vz= \vf(\vz, u)\frac{\dd u}{\dd t} \dd t + \vg(\vh,t)\dd t
\]
\[
\frac{\dd u}{\dd t} = \zeta(\vz,\vh)
\]
The choice of $\zeta$ gives control on how often a token is emitted. When solving this equation for $\Delta t = 1$, the Euler update rule becomes simple
\[
\vz_{t+1}= \vz_t + \vf(\vz_t,u)\zeta(\vz_t,u) + \vg(\vh,t)
\]
\[
u_{t+1} = u + \zeta(\vz_t,u)
\]
Given that a hidden state is computed for every $t$, we could use CTC
\[
\Ls = \mathbf{CTC}(P(\vy_{1:T}|\mZ_{1:T}),\vy^\mathbf{lbl}_{1:U})
\]
It is simple to add a rate penalty to steer towards correct emissions/spikes
\[
\Ls=\int_0^U\zeta(\vz)\dd t - U
\]


% have a bias to emit tokens depending on how many tokens left and 

\section{Validation results}

  \begin{table}[ht]
  \centering
  \small
  \begin{tabular}{lccccc}
  \hline
  Model & WER & GPUs & Steps & Train time & Peak GPU memory \\
  \hline
  Streaming AED
    & 14.89\% (AED) & 2 & 2{,}000 & $\sim$14 min
    & 0.53 / 0.62 GiB \\
  Compact MLP CTC
    & 19.86\% (frame CTC) & 2 & 4{,}000 & 0.41 h
    & 0.66 / 0.75 GiB \\
  Pretrained CTC head
    & 24.82\% (frame CTC) & -- & -- & --
    & -- \\
  Joint-clock ODE
    & 7.17\% (frame CTC)
    & 2 & 4{,}000 & 1.62 h & 1.37 / 1.77 GiB \\
    ODE adjoint
    & 9.93 & 1& 0.89 / 0.98 GiB \\
  \hline
  \end{tabular}
  \caption{VCTK validation results. GPU memory is peak allocated / reserved
  memory per GPU. The AED memory is estimated from a one-step replay and excludes
  DDP communication buffers. The pretrained CTC score is from a separate
  buffered-window diagnostic and is not directly matched to the other runs.}
  \label{tab:streaming-asr-results}
  \end{table}

  Encoder is frozen.